\section{Retained structure and
normalization}\label{retained-structure-and-normalization}

Write \(q=2^N\), where \(N>1\) is odd, and \(S_N=\operatorname{Sz}(q)\).
Use its usual ovoid action, with distinguished points \(\infty,0\), and
root group \(U_q\) at infinity. Write its central involutions as
\(z(b)\), \(b\in\mathbb F_q\), with \(z(b)z(c)=z(b+c)\). A Weyl
involution \(w\) swaps infinity and zero and is defined over
\(\mathbb F_2\).

We retain the following classical Suzuki-group facts. The ovoid and
root-group descriptions are the finite specialization of Lüneburg,
\emph{Translation planes} (1980), as cited in source Section 2.
The maximal-subgroup classification is Suzuki, \emph{On a class of doubly
transitive groups} (1962), cited explicitly in source Section 3.
The automorphism and centralizer facts belong to the same standard
structure theory; their use is retained, not a new classification here:

\begin{enumerate}
\def\labelenumi{\arabic{enumi}.}
\tightlist
\item
  The ovoid action is doubly transitive. Every involution fixes exactly
  one ovoid point. \(U_q\) acts regularly on the points other than
  infinity.
\item
  \(C_{S_N}(z(b))=U_q\) for \(b\ne0\), and all involutions in \(U_q\)
  are its central elements \(z(b)\). All involutions of \(S_N\) are
  conjugate.
\item
  The ordered two-point stabilizer is the split torus \(H_q\). Its
  conjugation acts as \(z(b)\mapsto z(tb)\) and
  \(wz(c)w\mapsto wz(t^{-1}c)w\), as \(t\) ranges bijectively over
  \(\mathbb F_q^*\). This uses a reparametrization of the usual diagonal
  torus. Its central-coordinate exponent is coprime to \(q-1\).
\item
  \(\operatorname{Aut}(S_N)=S_N\rtimes\langle F\rangle\), where
  \(F:x\mapsto x^2\) has order \(N\). Canonical subfield groups \(S_d\)
  for \(d\mid N\) have the corresponding subfield coordinates and
  ovoids.
\item
  Proper maximal subgroups are point stabilizers, odd-torus normalizers
  with cyclic Sylow 2-subgroups, and conjugates of proper subfield
  groups. A nonidentity odd-order element has a unique maximal odd
  cyclic torus as its centralizer. The normalizer of such a torus is
  proper.
\item
  \(S_N\) is nonabelian simple. String C-groups correspond to regular
  abstract polytopes, with reversing the generator order corresponding
  to duality.
\end{enumerate}

For clarity, the complete list of possible maximal-subgroup types used
here is the source's Table 1. Put $r=\sqrt{2q}$,
$\alpha=q+r+1$ and $\beta=q-r+1$. Then $|S_N|=q^2(q^2+1)(q-1)$ and
\[
\begin{array}{c|c}
\text{maximal subgroup type}&\text{order}\\\hline
 U_q\rtimes C_{q-1}&q^2(q-1)\\
 D_{2(q-1)}&2(q-1)\\
 C_{\alpha}\rtimes C_4&4\alpha\\
 C_{\beta}\rtimes C_4&4\beta\\
 \operatorname{Sz}(2^d),\quad N/d\text{ prime}&2^{2d}(2^{2d}+1)(2^d-1).
\end{array}
\]
Here $d$ is a proper divisor of $N$, and each index is $|S_N|$ divided
by the displayed order. The group $U_q$ has order $q^2$ and is the
source's nonsplit extension of two elementary abelian groups of order
$q$. Only this retained list and the stated involution and torus
properties enter the generation and intersection arguments.

The root/torus assertions can equally be read in the standard
four-by-four matrix model. Restricting them to a subfield is compatible
with the Suzuki field automorphism: since \(N/d\) is odd, its exponent
restricts to the canonical exponent for \(\mathbb F_{2^d}\). In
particular, the central torus scaling parameter lies in a subfield
exactly when its diagonal parameter does. The torus acts faithfully on
the nonzero central elements.

\textbf{Normalization theorem.} Inner-conjugacy classes of ordered
triples \((r_0,r_1,r_2)\) of involutions with \(r_0\ne r_2\),
\([r_0,r_2]=1\), and \([r_0,r_1]\ne1\) are in bijection with

\[X_N=(\mathbb F_q\setminus\mathbb F_2)\times\mathbb F_q^*.\]

The representative for \((b,c)\) is

\[\bigl(z(1),\ wz(c)w,\ z(b)\bigr).\]

\textbf{Proof.} Conjugate \(r_0\) to \(z(1)\). Its commuting involutions
are exactly \(z(b)\), and excluding the identity and \(r_0\) means
\(b\ne0,1\). Its centralizer \(U_q\) fixes every such \(z(b)\). The
involution \(r_1\) fixes a point other than infinity, since otherwise it
would belong to \(U_q\) and commute with \(r_0\). There is a unique
element of \(U_q\) that moves that point to zero. The resulting \(r_1\)
is a central involution of the root group at zero, hence \(wz(c)w\) with
\(c\ne0\). No further element of \(U_q\) fixes zero, proving uniqueness.
Conversely every displayed pair gives a triple with the stated
properties. \(\square\)

Thus field automorphisms act on normalized coordinates simultaneously:

\[F(b,c)=(b^2,c^2).\]

Reversing the endpoints gives \((z(b),wz(c)w,z(1))\). Normalize its
first entry by split-torus scaling \(t=b^{-1}\) to obtain the duality
action

\[D(b,c)=(b^{-1},bc).\]

It is an involution commuting with \(F\).

\section{Generation and the field of
definition}\label{generation-and-the-field-of-definition}

\textbf{Subfield-position lemma.} If a conjugate of a canonical \(S_d\),
\(d\mid N\), contains \(z(1)\) and \(wz(c)w\), then it is the canonical
\(S_d\) itself.

\textbf{Proof.} The embedded ovoid of that conjugate contains infinity
and zero, because those are the unique fixed ovoid points of the two
involutions. The canonical \(S_d\) is doubly transitive on its own
ovoid. Changing the conjugator on its \(S_d\) side therefore makes it
fix the ordered pair \((\infty,0)\), so it belongs to \(H_q\). Suppose
its central scaling is \(t\). The conjugate root center is then
\(z(t\mathbb F_{2^d})\). Since it contains \(z(1)\),
\(t\in\mathbb F_{2^d}^*\). The torus/subfield compatibility in Section 1
implies that the conjugator is in \(H_{2^d}\le S_d\). It consequently
normalizes \(S_d\), proving the claim. \(\square\)

\textbf{Generation theorem.} The group generated by the normalized
triple is exactly \(S_d\), where

\[\mathbb F_{2^d}=\mathbb F_2(b,c).\]

\textbf{Proof.} The displayed generators belong to \(S_d\), and \(d>1\)
because \(b\notin\mathbb F_2\). If they generated a proper subgroup of
\(S_d\), place it in a maximal subgroup. A point stabilizer cannot
contain the triple: its involutions would all fix that point, whereas
the first and middle generators fix different points. An odd-torus
normalizer cannot contain the four-group \(\langle z(1),z(b)\rangle\),
because its Sylow 2-subgroups are cyclic. The remaining possibility is a
conjugate of \(S_e\) for a proper divisor \(e\) of \(d\). The position
lemma makes it canonical, forcing \(b,c\in\mathbb F_{2^e}\), a
contradiction. \(\square\)

This proves generation from field coordinates without retaining the
source's enumeration formula or its generation-count lemma.

\section{The intersection property and
rank}\label{the-intersection-property-and-rank}

The first and middle involutions in a normalized triple do not commute;
neither do the middle and last, by their different fixed points. The
dihedral group generated by either adjacent pair has odd rotation order.
Indeed, if that order had 2-part at least four, its dihedral Sylow
2-subgroup would have noncommuting involutions, impossible inside a
Suzuki root group. If the rotation order had 2-part exactly two, its
central involution would centralize both generators, placing both in one
root group, whose involutions commute. Rotation order two is also
excluded by noncommutation.

Suppose these two odd-rotation dihedral subgroups intersect beyond their
common middle involution. Their intersection then contains a nontrivial
odd-order rotation \(x\). Both rotation subgroups lie in the unique
maximal torus \(C_{S_d}(x)\). Each adjacent dihedral subgroup normalizes
that torus, so all three generators lie in its proper normalizer. This
contradicts the generation theorem. Their intersection is consequently
exactly the order-two group generated by the middle involution.

The outer four-group intersects either odd-rotation dihedral subgroup in
exactly the appropriate endpoint subgroup, by order divisibility.
Together these statements give every rank-three subgroup-intersection
identity. Every normalized triple is therefore a nondegenerate string
C-group for its generated subfield group. Its adjacent product orders
are at least three.

There is no string C-group representation of \(S_N\) of rank at least
four. In such a representation, \(r_0\) commutes with every \(r_j\) for
\(j\ge2\). These involutions lie in the center of one root group and
therefore commute with each other. In particular \(r_3\) commutes with
all \(r_j\) for \(j\ge2\), as well as with \(r_0,r_1\) by the string
relations. This would be a nontrivial central element of \(S_N\). Ranks
one and two cannot generate a nonabelian simple group.

Rank three always occurs: take \(b\) of full field degree \(N\) and
\(c=1\). Finally, every involution of a group
\(S_N\le G\le\operatorname{Aut}(S_N)\) maps trivially to \(G/S_N\),
since that quotient has odd order. A proper extension cannot be
generated by involutions. This recovers the full rank/existence
statement of L085 quoted as L074, Theorem 3, without using that theorem
as an input.

\section{The two counting formulas}\label{the-two-counting-formulas}

Apply Theorem UD.ExactSupportOrbits to the divisor poset of \(N\). The
normalization and generation theorems prove that the support is the
joint field degree. If \(A(d)\) counts data defined over
\(\mathbb F_{2^d}\), the exact-degree count is \[
E(N)=\sum_{d\mid N}\mu(N/d)A(d).
\tag{E3}
\] Every exact-degree datum has a Frobenius orbit of length \(N\). The
additional duality quotient requires the freeness argument that follows;
it is not an automatic division by two.

For pairs of distinct commuting involutions, normalized data are just
\(b\in\mathbb F_q\setminus\mathbb F_2\). Put

\[\lambda(d)=\frac1d\sum_{e\mid d}\mu(d/e)2^e.\]

For \(d>1\), this is the number of Frobenius orbits of elements of exact
degree \(d\); the two degree-one elements have already been excluded.
Unordered-pair reversal sends \(b\) to \(b^{-1}\). Inversion commutes
with Frobenius and has no fixed point on these data: \(b^2=1\) would
give \(b=1\) in characteristic two. The acting Frobenius group has odd
order \(N\), so UD.OddOrbits proves that reversal fixes no Frobenius
orbit. Therefore the number of unordered-pair classes is

\[\boxed{\frac12\sum_{d\mid N,\ d>1}\lambda(d).}\]

This is exactly L074, Theorem 1.

For polyhedra whose full group is \(S_N\), the marked data are \(X_N\)
with exact joint field degree \(N\). Here

\[A(d)=(2^d-2)(2^d-1).\]

Their Frobenius orbits all have length \(N\). The normalized duality
\(D(b,c)=(b^{-1},bc)\) is an involution commuting with Frobenius, as
proved in Section 1. It has no fixed point because its first coordinate
would force \(b=1\). It preserves the joint field, since each of
\(b,c\) can be recovered from \(b^{-1},bc\). Thus UD.OddOrbits applies
on the exact-support subset, and the free-action consequence of
UD.AdmissibleExtensions gives the number of polyhedra up to isomorphism
and duality:

\[\boxed{P(N)=\frac1{2N}\sum_{d\mid N}
   \mu(N/d)(2^d-2)(2^d-1).}\]

There are no self-dual examples. The formula also gives the complete
polytope count, by the rank restriction in Section 3. For prime \(N\) it
reduces to \((q-2)(q-1)/(2N)\), the source's Theorem 4. For
\(N=3,5,9,15\) it gives respectively \(7,93,14476,35788085\).

\section{Equivalence with the source's nested
formula}\label{equivalence-with-the-sources-nested-formula}

This section proves equivalence for every odd \(N\), not just agreement
of a finite table. Let \(Q(N)\) be the number of Frobenius orbits on all
of \(X_N\), including data of smaller joint degree.

Apply UD.ExtensionOrbits to the projection onto the first coordinate,
and stratify by its degree \(n>1\). There are \(\lambda(n)\)
first-coordinate orbits. After choosing a representative, its stabilizer
is generated by \(F^n\) and acts on \(\mathbb F_q^*\). For
\(m\mid N/n\), the number of points fixed by \((F^n)^m\) is
\(2^{nm}-1\). Möbius inversion gives

\[\psi(n,N)=\sum_{m\mid N/n}\frac1m
       \sum_{e\mid m}\mu(m/e)(2^{ne}-1)\]

orbits of this stabilizer. Therefore

\[Q(N)=\sum_{n\mid N,\ n>1}\lambda(n)\psi(n,N).\]

Alternatively stratify by the joint field degree \(d\). Orbits on
exact-degree-\(d\) pairs have length \(d\), and their Frobenius
equivalence is the same whether computed in \(\mathbb F_{2^d}\) or the
larger field. Thus \(Q(N)=\sum_{d\mid N}E(d)/d\). Möbius inversion now
gives

\[\frac12\sum_{d\mid N}\mu(N/d)Q(d)=\frac{E(N)}{2N}=P(N).\]

The expression on the left is precisely L074, Theorem 2. This derives
its full parameter range and nondegeneracy conclusion from the marked
field theorem, with all field-automorphism orbit corrections included.

\section{Source statements and retained
background}\label{source-statements-and-retained-background}

The preceding proof also locates the source's auxiliary counting
statements. Lemma 7 is the divisor-poset instance of Möbius inversion
proved in UD.ExactSupportOrbits; Lemma 8 is the formula for \(\lambda\);
Lemma 9 is the stabilizer-fiber formula for \(\psi\); and Lemma 10 is
\(Q(N)/2\), before removing data of smaller joint field degree. Remark 1
is the prime-degree specialization. The source's second example follows
by substituting \(N=15\): the values of \(\lambda\) at \(3,5,15\) are
\(2,6,2182\), the corresponding values of \(\psi(n,15)\) are
\(6559,10943,32767\), and \(Q(15)/2=35788185\). Removing the
smaller-degree counts \(7\) and \(93\) gives \(35788085\). Its first
example illustrates the same stabilizer \(\langle F^n\rangle\): a
degree-three coordinate in \(\mathbb F_{2^{15}}\) is fixed by \(F^3\),
whereas a general second coordinate is not.

The earlier structural lemmas have a different status. The source
attributes Lemmas 1--3 to Lüneburg (1980), Lemma 4 to Leemans (1998),
and Table 1 to Suzuki (1962); those attributions are retained here. Source Lemma 1
follows from the cyclic automorphism group of a finite field: a
Frobenius power \(F^j\) has square \(F\) exactly when
\(2j\equiv1\pmod N\), which has a unique solution for odd \(N\) and none
for even \(N\). The finite ovoid action and root-group consequences of
Lemma 2, the ovoid size in Lemma 3, and the maximal-subgroup
classification in Table 1 are retained in Section 1. We do not rederive
Lemma 2's arbitrary-field Bruhat description or Lemma 3's assertion
about plane sections of general ovoids. Those broader statements are not
premises of our counting proof. Lemma 4's torus-order coprimality is
likewise not used in place of the intersection argument: Section 3 uses
uniqueness of the centralizer torus and proves the intersection
directly. Lemma 5's involution count is \((q^2+1)(q-1)\) because each of
the \(q^2+1\) ovoid points supports \(q-1\) central root involutions and
every involution fixes exactly one point; its conjugacy assertion is
among the retained structure facts. Lemma 6 is the classical Fermat
congruence, and the definitions of ovoids, Möbius inversion and the
counting function are inherited notation. This inventory distinguishes
recovery of the principal classification from reproving every background
statement printed in the source.
